problem:  
Let \((M^{4n}, g)\) be a compact, simply connected, positive quaternion‑Kähler manifold with \(\operatorname{Hol}(g)\subset \mathrm{Sp}(n)\mathrm{Sp}(1)\), \(n\ge 2\).  
Denote the curvature operator \(\mathcal R_p:\Lambda^2(T_pM)\to \Lambda^2(T_pM)\), and let the standard \(\mathrm{Sp}(n)\mathrm{Sp}(1)\)-holonomy decomposition of the 2‑forms be  
\[
\Lambda^2(T_pM)= \mathfrak{sp}(n)\;\oplus\; \mathfrak{sp}(1)\;\oplus\;\mathcal W_p,
\]
where \(\mathcal W_p\) corresponds to the quaternionic Weyl component.  
For each point \(p\in M\), restrict \(\mathcal R_p\) to \(\mathcal W_p\) and denote its scalar invariants by the ordered tuple of eigenvalues (with multiplicity)
\[
\sigma(p):=\text{spec}\big(\mathcal R^{\mathcal W}_p\big)\subset \mathbb{R}^k.
\]
It is known that for any Wolf space (compact symmetric quaternion‑Kähler manifold) this spectrum is constant over \(M\).  

**Problem:**  
Determine whether the converse holds: If \(\sigma(p)\) is constant on \(M\), must \((M,g)\) be locally symmetric?  
Equivalently, does constancy of the Weyl curvature operator spectrum on the holonomy‑irreducible component \(\mathcal W\) imply \(\nabla R=0\) for a compact positive quaternion‑Kähler manifold?  

Why is it a "good" problem:  
This problem is conceptually precise, geometrically natural, and analytically nontrivial. It asks whether purely spectral data encoded in the curvature operator suffice to detect local symmetry within the deeply constrained quaternion‑Kähler category. The condition is weaker than the standard requirement \(\nabla R=0\) yet stronger than known scalar identities, making it a subtle rigidity question that intertwines geometric analysis, representation theory of \(\mathrm{Sp}(n)\mathrm{Sp}(1)\), and holonomy theory.

A positive resolution would yield a new spectral curvature characterization of the Wolf spaces among all positive quaternion‑Kähler manifolds, providing a curvature‑analytic route toward the LeBrun–Salamon conjecture. Even a counterexample would reveal new phenomena in the interplay between curvature spectra and special holonomy geometry. The statement is simple to express—“Does constant curvature spectrum force symmetry?”—but its depth is considerable, uniting geometry, spectral theory, and global symmetry in a single, elegant question.